\begin{itemize}
\item Imposed an upperbound of 14 kWDC, which contains 95\% of all installations. Counties with source\_count <10 are backfilled with aggregates at the state level. Distribution based on all installations is applied only to occupied single-family detached  homes, actual distribution for single-family detached homes may be higher.
\item PV is not modeled in AK. No data has been identified.
\item For Hawaii, EIA Form 861 reports the number of residential utility customers with rooftop PV in the Net Metering 2023 worksheet and the total number of residential utility customers in Sales Ult Cust 2023 worksheet. Taking the ratio of these two gives us the fraction of homes with rooftop PV. We place the additional constraint that all of these rooftop PV customers are in single-family detached homes. Because Hawaii is composed of multiple islands, the four electric utilities in Hawaii (and which report to EIA Form 861) approximately align with the county boundaries with Hawaii, Honolulu, and Kauai having one county for one electric utility, and both Maui and Kalawao mapping to the same electric utility, so we we develop distinct PV saturation rates for each county in Hawaii (note Kalawao has population 81 people and is located on the island of Maui).
\end{itemize}
